% Chapter 3: Single-User Proactive Delivery (v2 ready-aware)
% Data source: paper_materials/proactive_single_user_v2/runs/2026-05-24T16-06-30
% Delivery: HPC origin (10.120.17.176:8090), interleaved ordering, N=3 per baseline
% Generated: 2026-05-25

\section{Single-User Proactive Delivery}
\label{sec:single-proactive}

\subsection{Motivation}

On-demand HTTP Range fetching of large 3D Gaussian Splatting (3DGS) tile files
introduces a fundamental timing problem: the client cannot render a viewport
region until the corresponding tile chunk has been fetched, decoded, and
uploaded to the GPU. For a 1.28\,GB \texttt{.rad} file served over a campus
LAN (round-trip time $\approx$8\,ms, effective throughput $\approx$80\,MB/s),
each 1.7\,MB chunk requires $\approx$80--120\,ms from request to render-ready
state. When the viewer's camera moves into a new region, the on-demand
pipeline cannot meet a 100\,ms interactive deadline for the majority of
demanded chunks.

We quantify this problem and evaluate whether \emph{ready-aware proactive
delivery}---prefetching and pre-parsing chunks before they are demanded by the
renderer---can reduce deadline miss rates without degrading render performance.

\subsection{Experimental Design}

\paragraph{Asset and delivery.}
We use a single 40-million-splat scene (Coit Tower, San Francisco) encoded as
a chunked \texttt{.rad} file (1.28\,GB, 341 chunks at LOD-0). The file is
served from an HPC origin server over campus LAN via HTTP Range requests
(verified HTTP/1.1 206 Partial Content with \texttt{Accept-Ranges: bytes}).

\paragraph{Baselines.}
Three delivery strategies are compared under identical scene, camera trace,
and measurement protocol:
\begin{itemize}
  \item \textbf{B0 (On-demand)}: Spark's default streaming-LOD pipeline.
    Chunks are fetched only when the pager determines they are needed for the
    current viewport.
  \item \textbf{B3-ready (VRC-Single v2)}: A ready-aware proactive controller
    that predicts camera pose 0.5--2\,s ahead, scores chunks by viewport
    relevance and LOD importance, and prefetches the top-ranked chunks via a
    Web Worker with bounded concurrency ($\text{maxParallel}=2$, byte budget
    128\,MB). On fetch completion, the controller emits a
    \texttt{parse\_complete} event so the chunk's ready time advances to
    parse-complete rather than remaining at network-finish.
  \item \textbf{B4-ready-oracle}: Same ready-aware pipeline, but chunk
    priority is determined by perfect future knowledge of the B0 demand trace
    (oracle upper bound on prediction accuracy).
\end{itemize}

\paragraph{Measurement protocol.}
Each trial consists of a 20\,s warmup phase followed by a 40\,s measurement
phase with automated camera orbit. Trials are run in \emph{interleaved} order
(B0/B3/B4/B0/B3/B4/B0/B3/B4) so that each baseline experiences the same
progressive cache-warming sequence. All trials use a local instrumented proxy
server that logs per-request server-side timings aligned with client-side CDP
network audit rows (alignment rate $\geq$99.7\%).

\paragraph{Metrics.}
\begin{itemize}
  \item \texttt{deadline\_miss\_ratio\_100ms}: fraction of demanded chunks
    whose ready time exceeds needed time by more than 100\,ms.
  \item \texttt{useful\_chunks\_before\_demand}: count of chunks that reach
    ready state before the renderer demands them.
  \item \texttt{first\_visible\_splat\_ms}: time from page load to first
    non-zero visible splat count.
  \item \texttt{wasted\_prefetch\_ratio}: bytes prefetched but never demanded,
    divided by total prefetched bytes.
  \item \texttt{measure\_fps}: mean render frame rate during the measurement
    phase.
\end{itemize}

\subsection{Results}

Table~\ref{tab:v2-aggregate} summarizes the aggregate results across $N=3$
trials per baseline on the HPC origin delivery.

\begin{table}[t]
\centering
\caption{Single-user proactive delivery: aggregate metrics (mean $\pm$ s.d.,
$N=3$ per baseline, HPC origin, interleaved trial order).}
\label{tab:v2-aggregate}
\begin{tabular}{lcccccc}
\hline
Baseline & miss$_{100}$ (\%) & useful & first-vis (s) & FPS & wasted (\%) & demanded \\
\hline
B0 On-demand      & $69.8 \pm 26.3$ & $106 \pm 94$ & $5.0 \pm 1.9$ & $107.4 \pm 34.5$ & $0.0$ & $305 \pm 89$ \\
B3-ready          & $62.5 \pm 5.2$  & $121 \pm 13$ & $3.7 \pm 0.7$ & $88.5 \pm 3.0$   & $6.2 \pm 0.8$ & $324 \pm 15$ \\
B4-ready-oracle   & $66.5 \pm 3.7$  & $104 \pm 12$ & $7.7 \pm 1.3$ & $86.8 \pm 1.3$   & $10.5 \pm 0.4$ & $312 \pm 4$ \\
\hline
\end{tabular}
\end{table}

\paragraph{Cold-start performance.}
The most informative comparison is the first trial of each baseline (t0),
which represents a user's initial visit to the scene with no cached data
(Table~\ref{tab:v2-coldstart}).

\begin{table}[t]
\centering
\caption{Cold-start trial (t0): per-trial metrics with no prior cache state.}
\label{tab:v2-coldstart}
\begin{tabular}{lccccc}
\hline
Baseline & miss$_{100}$ (\%) & useful & first-vis (ms) & FPS & demanded \\
\hline
B0 On-demand    & 99.5 & 0   & 7{,}111 & 147.2 & 203 \\
B3-ready        & 65.1 & 119 & 3{,}352 & 85.5  & 341 \\
B4-ready-oracle & 69.7 & 94  & 8{,}500 & 87.2  & 310 \\
\hline
\end{tabular}
\end{table}

In the cold-start condition, B3-ready reduces deadline miss by 34.4
percentage points ($-34.6\%$ relative to B0) and halves first-visible latency
from 7.1\,s to 3.4\,s ($-52.9\%$). The controller enables Spark's pager to
demand 341 chunks within the measurement window (vs.\ 203 for B0), because
prefetched chunks complete faster and unblock subsequent paging decisions.

\paragraph{Cache-warming dynamics.}
Figure~\ref{fig:cold-warm}d--e shows how B0's performance improves
dramatically across successive trials as the operating system and browser
caches warm: miss$_{100}$ drops from 99.5\% (t0) to 49.6\% (t2), and
useful-before rises from 0 to 178. B3-ready remains stable across rounds
(miss$_{100}$ = 65.1\%, 56.6\%, 65.9\%), because its prefetch schedule does
not depend on host cache state. This reveals that B0's aggregate mean
(69.8\%) is dominated by warm-cache trials where Spark's natural pipeline
achieves near-equivalent pre-staging without the proactive controller.

\paragraph{First-visible latency.}
B3-ready achieves a mean first-visible of 3{,}746\,ms vs.\ B0's 4{,}988\,ms
($-24.9\%$; Fig.~\ref{fig:first-visible}b). The improvement is concentrated
in the cold-start trial (B3: 3{,}352\,ms vs.\ B0: 7{,}111\,ms). In warm
trials, both baselines converge to $\approx$3.5--4.5\,s.

\paragraph{Oracle paradox.}
Counterintuitively, B4-ready-oracle produces \emph{worse} first-visible
latency (7{,}673\,ms mean) than B3-ready (3{,}746\,ms). The oracle sorts
prefetch priority by the B0 reference demand trace, which places the .rad
metadata header (chunk index 0, 65\,KB) first. This delays viewport-relevant
chunks that determine first-visible. B3-ready's viewport-priority scoring
naturally prioritizes spatially relevant chunks, producing faster initial
content appearance despite imperfect prediction.

\paragraph{Prefetch efficiency.}
B3-ready wastes only 6.2\% of prefetched bytes on chunks never demanded
(Fig.~\ref{fig:efficiency}g), well below the 30\% threshold. The
viewport-priority scoring function effectively targets chunks the renderer
will need within the prediction horizon.

\paragraph{Render performance.}
B3-ready maintains 88.5\,FPS (vs.\ B0's 107.4\,FPS; $-17.6\%$), remaining
above the 60\,FPS interactive threshold. The FPS difference is partly an
artifact of B0's inflated t0 value (147.2\,FPS with only 203 demanded chunks,
indicating the renderer was idle waiting for data rather than rendering at
full load). In warm trials, B0 and B3 converge to $\approx$87--91\,FPS.

\subsection{Discussion}

The results demonstrate that ready-aware proactive delivery provides a
measurable benefit in the cold-start regime---the scenario most relevant to a
new user joining a 3DGS scene. The controller's value is concentrated in the
first $\approx$5--10 seconds of a session, during which on-demand fetching
leaves the renderer starved for data. Once host-side caches warm (after the
first full traversal of the scene), Spark's natural streaming-LOD pipeline
achieves comparable pre-staging rates without the proactive controller.

This finding has two implications for multi-user systems:
\begin{enumerate}
  \item For the \emph{first user} in a session (or any user visiting a
    previously uncached region), proactive delivery provides substantial
    cold-start acceleration.
  \item For \emph{subsequent users} whose viewports overlap with prior
    visitors, the host-side cache warming effect suggests that multi-user
    coordination could amplify the benefit by ensuring that one user's
    prefetch warms the cache for others---motivating the viewport-aware range
    coordination (VRC) studied in subsequent chapters.
\end{enumerate}

\subsection{Limitations}

\begin{itemize}
  \item The ready-event policy uses \texttt{parse\_complete} as a proxy for
    GPU-upload-complete, because Spark's current architecture does not expose
    per-chunk upload events. The actual render-ready time may be slightly later
    than reported.
  \item The interleaved trial ordering controls for cache warming but does not
    eliminate it. A fully randomized design with explicit cache flushes between
    trials would provide tighter confidence intervals.
  \item The measurement is conducted on a single scene (Coit Tower, 40M
    splats). Generalization to other scene sizes and network conditions
    requires additional experiments.
\end{itemize}

\begin{figure}[t]
  \centering
  \includegraphics[width=\linewidth]{figures/proactive_single_user_v2/fig_v2_deadline_miss.pdf}
  \caption{\textbf{Deadline miss ratio} (ready time exceeds needed time by
  $>$100\,ms). Individual trial values shown as dots; bars indicate mean
  $\pm$ s.d. ($N=3$). B3-ready shows lower variance and lower mean miss rate
  than B0, whose high variance reflects progressive cache-warming across trials.}
  \label{fig:deadline-miss}
\end{figure}

\begin{figure}[t]
  \centering
  \includegraphics[width=\linewidth]{figures/proactive_single_user_v2/fig_v2_first_visible.pdf}
  \caption{\textbf{First visible splat latency.} B3-ready achieves
  3.7\,s mean first-visible vs.\ B0's 5.0\,s ($-24.9\%$). B4-oracle's
  poor first-visible (7.7\,s) results from demand-trace-ordered prefetch
  that delays viewport-relevant chunks.}
  \label{fig:first-visible}
\end{figure}

\begin{figure}[t]
  \centering
  \includegraphics[width=\linewidth]{figures/proactive_single_user_v2/fig_v2_cold_warm_dynamics.pdf}
  \caption{\textbf{Cold-start vs.\ warm-state dynamics.} (d)~Deadline miss
  per trial round: B0 drops from 99.5\% to 49.6\% as caches warm; B3-ready
  remains stable at $\approx$60--66\%. (e)~Useful chunks before demand: B0
  rises from 0 to 178 with warming; B3-ready is consistently 109--134.}
  \label{fig:cold-warm}
\end{figure}

\begin{figure}[t]
  \centering
  \includegraphics[width=\linewidth]{figures/proactive_single_user_v2/fig_v2_efficiency.pdf}
  \caption{\textbf{Efficiency metrics.} (f)~Render FPS: all baselines exceed
  60\,FPS; B0's elevated t0 value (147\,FPS) reflects renderer idling while
  waiting for data. (g)~Wasted prefetch: B3-ready wastes only 6.2\% of
  prefetched bytes, well below the 30\% threshold (red dashed line).}
  \label{fig:efficiency}
\end{figure}

\begin{figure}[t]
  \centering
  \includegraphics[width=\linewidth]{figures/proactive_single_user_v2/fig_v2_visibility_timeline.pdf}
  \caption{\textbf{Visibility timeline} (cold-start trial t0). B3-ready
  achieves non-zero visible splats $\approx$3.8\,s earlier than B0, and
  reaches higher splat counts faster due to prefetch-enabled paging.}
  \label{fig:visibility-timeline}
\end{figure}
